\documentclass[10pt]{article}
\usepackage[margin=1in]{geometry}
\usepackage{booktabs,graphicx,hyperref,amsmath}
\usepackage[numbers]{natbib}
\title{VITRINE: Explainable Static PE Triage and Attributable Temporal Drift\\from EMBER 2018 to EMBER2024}
\author{Rakshit\\VIT Chennai}
\date{Preprint, October 2026 (not peer reviewed)}
\begin{document}
\maketitle

\begin{abstract}
Static machine-learning malware detectors are usually reported as a single score on a single,
temporally adjacent test set. We present VITRINE, a static-only triage pipeline that scores Windows PE
files with gradient-boosted trees over 91 named features, explains each verdict with exact TreeSHAP,
tags import-level capabilities and synthesizes benign-validated YARA rules. Its named features let us
measure \emph{and attribute} temporal drift. Trained on EMBER 2018 and tested on EMBER2024, the model
loses 0.018 ROC AUC (0.9914 to 0.9730, 95\% CI 0.9710--0.9748) and, at its own calibrated threshold,
recall falls from 88.0\% to 58.6\% while the false-positive rate \emph{drops} to 0.07\%: the decay is
silent. A 2024 model scored on 2018 data fails loudly instead (44.8\% FPR). The shift concentrates in
a few named artifacts (embedded executables, DLL share, signing, toolchain versions), but removing them
does not recover robustness. We also report honestly where VITRINE is behind: on identical test rows the
EMBER authors' published 2018 model is ahead by 0.0047 AUC and 12.7 points of detection at 0.1\% FPR.
All experiments use pre-extracted features only; no binaries are handled.
\end{abstract}

\section{Introduction}
Analysts need more than \texttt{predict() = 0.87}: they need reasons, a candidate signature, and an
idea of whether the model still works on today's files. EMBER \citep{anderson2018ember} established
static-feature baselines; EMBER2024 \citep{joyce2025ember2024} renewed the data; TESSERACT
\citep{pendlebury2019tesseract} showed that evaluation must respect time. We combine these: an
interpretable triage model whose features have names, evaluated across two dataset generations six
years apart, with drift attributed to those names. Contributions: (i) a cross-time evaluation
EMBER 2018 $\leftrightarrow$ EMBER2024 with seed-pooled bootstrap intervals, a same-size control and a
schema-translation control; (ii) a feature-level drift attribution and a negative ablation result;
(iii) a re-scoring of both datasets' published reference models on our exact rows, so that the cost
of interpretability is measured against the real state of the art rather than an under-resourced
reproduction; (iv) a benign-validated structural YARA synthesizer with abstention, compared with
like-for-like benign-filtered baselines.

\section{Related work}
Static PE detectors with hashed features and GBDTs \citep{anderson2018ember,ke2017lightgbm} are strong
but opaque. SHAP \citep{lundberg2017shap} and exact TreeSHAP \citep{lundberg2020tree} give additive
attributions for tree ensembles \citep{chen2016xgboost}. Automatic rule generation ranges from
frequency-based tools such as yarGen \citep{yargen} to biclustering \citep{raff2020autoyara}.
capa \citep{capa} detects capabilities from code. Time-aware evaluation \citep{pendlebury2019tesseract}
motivates our cross-period design.

\section{Method}
\textbf{Parsing and features.} A bounds-checked pure-Python PE parser emits the EMBER v2 raw schema;
91 named features (header flags, mitigations, section entropy, import groups, string counters,
capability indicators) are computed from it. EMBER2024 feature-version-3 records are translated to v2
(section names lower-cased, a few counters approximated).
\textbf{Model and explanation.} XGBoost (600 trees, depth 10); thresholds at 1\% and 0.1\% FPR
calibrated on a held-out period; exact TreeSHAP contributions with sentence templates.
\textbf{Rules.} Structural atoms (imports, section names, imphash) kept only if rare in a training
benign pool; the smallest $N$-of threshold with zero benign hits, else abstain.
\textbf{Drift attribution.} Per feature, population stability index between periods (per class)
times the 2018 model's mean $|$SHAP$|$; ablation retrains without the top-10 features.

\section{Experimental setup}
EMBER 2018: deterministic 50\% sha256-prefix subsample of the 600k labelled training rows (memory
limit of a 16\,GB laptop), 275{,}732 rows to fit, 2018-10 for calibration, 200{,}000 test files.
EMBER2024 Win32: first 8\,MiB of each of 64 weekly files (a sha256-range sample of about 7\%;
135{,}423 files), weeks 0--47 fit, 48--51 calibrate, 52--63 test. Three seeds per trained model;
1{,}000-resample stratified bootstrap pooled over seeds; paired differences on shared resamples.

\section{Results}
\begin{table}[h]\centering\small
\begin{tabular}{lcccc}\toprule
Train $\rightarrow$ test & ROC AUC (95\% CI) & TPR@1\% & FPR / recall at source thr.\\\midrule
2018 $\rightarrow$ 2018 & 0.9914 (0.9908--0.9922) & 88.2\% & 0.93\% / 88.0\%\\
2018 $\rightarrow$ 2024 & 0.9730 (0.9710--0.9748) & 74.2\% & 0.07\% / 58.6\%\\
2018 (99k) $\rightarrow$ 2024 & 0.9712 (0.9689--0.9738) & 73.8\% & 0.02\% / 51.7\%\\
2024 (99k) $\rightarrow$ 2024 & 0.9946 (0.9940--0.9951) & 91.1\% & 1.18\% / 91.8\%\\
2024 (99k) $\rightarrow$ 2018 & 0.8860 (0.8810--0.8895) & 33.3\% & 44.8\% / 95.5\%\\\bottomrule
\end{tabular}
\caption{Cross-time results, 91 named features, XGBoost, 3 seeds.}\end{table}

At equal training size, an in-period 2024 model beats the 2018 model on 2024 by +0.023 AUC
(+0.021, +0.026). The 2018 model's weekly AUC over EMBER2024 ranges 0.918--0.992. The schema
translation itself costs only 0.0006 AUC (0.0003--0.0008) when the EMBER2024 paper configuration is
trained on native vs translated vectors of identical rows. That control shows that little information
is lost inside 2024; it cannot detect a feature defined differently in 2018 and 2024, and the drop
still mixes time with collection and extractor changes.

\textbf{Attribution.} The raw top drift-weighted feature, \texttt{n\_embedded\_mz}, is a definition
change: EMBER v2 counts the bytes \texttt{MZ} anywhere in a file, the v3 record only strings containing
\texttt{!This program}. On the same 600 benign System32 files the two definitions agree exactly on 13\%
of files and differ by PSI 8.9, more than its cross-dataset PSI (3.7). We therefore exclude 4 redefined
string counters and 6 proxy flags. Among the remaining features, \texttt{is\_dll} and \texttt{cert\_kb}
shift in benign files only (PSI$_\text{mal}\le0.003$), pointing to benign-sample selection, while
\texttt{n\_imports}, \texttt{os\_major}, \texttt{size\_kb}, \texttt{linker\_major}, \texttt{code\_ratio} and
\texttt{avg\_string\_len} shift in both classes; these are still computed by different parsers (LIEF vs
pefile). Retraining the 2018 model (99k rows, 3 seeds) without the top ten faithful features changes
2024 AUC by $-0.0040$ ($-0.0059$, $-0.0016$) at a 2018 cost of $-0.0025$; 20 group-matched random
draws of ten faithful features average $-0.0011$ ($-0.0035$ to $+0.0015$) and none is worse in AUC; at TPR@0.1\% FPR, 14 of 20 random draws are worse than the top ten
(mean $-0.063$ vs $-0.043$). The ranking is a correlational hypothesis about the source of drift, not an
established attribution, and it is not removable by deleting the named features.

\textbf{EMBER paper reproduction.} Under the published setup (EMBER 2017 feature v1, 2{,}351 dims,
elastic/ember vectorizer, LightGBM defaults, 600k labelled rows, run in GitHub Actions) we obtain AUC
0.99911 (0.99905--0.99917), 98.2\% and 93.0\% TPR at 1\% / 0.1\% FPR, against the paper's 0.99911,
98.2\% and 92.99\%; all paper values lie inside our intervals.

\textbf{Reference models.} The EMBER authors' published 2018 model, scored by us on the same 200k test
rows, reaches AUC 0.9964, 96.5\% / 87.0\% TPR at 1\% / 0.1\% FPR (consistent with their notebook, 86.8\% at 0.1\% FPR, within our CI; scored on our clean-room vectors, whose 50 entry-section-name hash dimensions differ from elastic/ember); VITRINE
is behind by 0.0047 AUC (0.0045--0.0049), 7.8 and 12.7 points. The released EMBER2024 Win32 model scores
AUC 0.9983 on our test subsample (paper: 0.9984).

\textbf{Rules.} Over 15 families out of time, VITRINE's rules cover 16.8\% of test siblings with 5
false positives per 100k benign files, against 13.9\% and 7 for an imphash baseline filtered on the
same benign pool; four of VITRINE's eight rules barely generalise.

\textbf{Robustness.} In a feature-space attack combining overlay, section, import and header edits,
detection falls to about 40--44\% for all models; an import-based capability floor appears to help,
but most of that gain is triggered by the benign donor's own imports, and it costs 11.7\% benign FPR.

\section{Limitations}
EMBER 2018 training uses half of the labelled rows, so the gap to the published model mixes data and
representation. EMBER2024 is subsampled. 
Adversarial edits are simulated in feature space. Capability tags were compared with capa only as
agreement, not ground truth.

\section{Ethics and reproducibility}
No malware binaries were downloaded, stored or executed; all data are pre-extracted features
(EMBER 2018: MIT; EMBER2024: Apache-2.0) plus read-only scans of local benign Windows files. Code,
result files, per-range SHA-256 manifests and exact commands are at
\url{https://github.com/rakshit-737/vitrine}.

\section*{Acknowledgements}
Code, experiments and this text were developed with AI assistance (Claude, Anthropic); the author
reviewed and is responsible for the content.

\bibliographystyle{plainnat}
\bibliography{refs}
\end{document}
